\documentclass[11pt]{article}

\usepackage[utf8]{inputenc}
\usepackage[margin=0.85in]{geometry}
\usepackage{amsmath,amssymb,mathrsfs,mathtools}
\usepackage{tikz-cd}
\usetikzlibrary{decorations.pathmorphing}
\usepackage{booktabs,array}
\usepackage{enumitem}
\usepackage{xcolor}
\usepackage[colorlinks=true,linkcolor=blue,urlcolor=blue]{hyperref}

\newcommand{\OO}{\mathbb O}
\newcommand{\CC}{\mathbb C}
\newcommand{\RR}{\mathbb R}
\newcommand{\Ostar}{\OO^{*}}
\newcommand{\cA}{\mathscr A_A}
\newcommand{\TA}{\mathcal T_A}
\newcommand{\In}{\operatorname{In}}
\newcommand{\Out}{\operatorname{Out}}
\newcommand{\Grpd}{\mathbf{Grpd}}

\definecolor{deepblue}{RGB}{31,78,121}
\definecolor{forest}{RGB}{35,105,66}
\definecolor{warmgray}{RGB}{85,85,85}

\title{\textbf{The Octonionic Round Trip and Its Genuine Total}\\
\large Insertion outline for the Concentricity proof}
\author{Jesse Michael Paul}
\date{24 July 2026}

\begin{document}
\maketitle

\noindent\textcolor{warmgray}{\textbf{Purpose.}}
This is a structural insertion outline, not final exposition.  It belongs
after the slice-wise construction of the distinguished
Euler--Weierstrass--GPV action and before the application of Riehl 8.3.4,
$\pi_0$, and $\operatorname{val}$.  Its purpose is to make explicit the
arrow-level passage from the slice geometry to the genuine total of the
octonionic round trip.

\section{The three categorical levels}

\subsection{The slice projection}

The slice-wise geometric action is the functor
\[
  A_{\mathrm{slice}}\colon
  \mathrm{GreatCircle.Base}\longrightarrow\mathrm{SphereWorld}.
\]
It contains the two-point stabilizer, the unique orbit--stabilizer
factorization, the full prime-indexed Euler lift, its GPV winding data, and
the pointwise tame continuous triangles to $N$.  It is essential, but it is
the sphere/Moebius projection of the action rather than the octonionic
round trip itself.

\subsection{The octonionic round-trip functor}

The native home of the A-section is $\Ostar$.  Its actual round trip is the
$G_2$-equivariant endofunctor
\[
  A_{\OO}\colon H_1\longrightarrow H_1,
  \qquad s\longmapsto A.\operatorname{realize}(s).
\]
The normalized state world has natural input and output functors, and the
round trip is their factorization:
\[
\begin{tikzcd}[column sep=large,row sep=large]
  \mathrm{StateWorld}_A
    \arrow[r,"\In"]
    \arrow[dr,"\Out"']
  & H_1 \arrow[d,"A_{\OO}"] \\
  & H_1
\end{tikzcd}
\qquad
\Out=A_{\OO}\circ\In.
\]
This triangle is arrow-level data.  It says that the input and output move
naturally under the same $G_2$ direction transport; it is stronger than the
object assignment $s\mapsto A.\operatorname{realize}(s)$.

\subsection{The projective indexing diagram}

The complete normalized state and its analytic presentation are then
curried over the projective base:
\[
  \cA\colon \mathrm{GreatCircle.Base}\longrightarrow\Grpd.
\]
For a projective object $X$,
\[
  \cA(X)
  =
  \mathrm{StateWorld}_A
  \times
  \mathrm{PresentationWorld}_A(X).
\]
The first factor contains the physical input/output round trip.  The second
contains the full Euler--Weierstrass--GPV tape and all of its pointwise
triangles to $N$.  Thus $\cA$ is the Grothendieck-indexing diagram of the
round-trip geometry, not a second octonionic endofunctor.

\section{The distinguished element and the two sweeps}

\subsection{Vertical: the entire prime stack through $N$}

For every Euler path $\delta$ the lift is literally
\[
  L_\delta(t)
  =
  \sum\nolimits'_{p:A.\iota} A.\ell_p(\delta(t)).
\]
The prime index remains internal to the infinite sum.  At every instant,
\[
\begin{tikzcd}[column sep=huge,row sep=large]
  L_\delta(t)
    \arrow[r,"\exp"]
    \arrow[d,rightsquigarrow,
      "\text{unique tame lift}"']
  & A.F(\delta(t))
    \arrow[d,"\text{diagonal Moebius action}"] \\
  N\text{-positioned lift}
    \arrow[r,"\text{same action at }N"']
  & N\text{-positioned value}
\end{tikzcd}
\]
commutes.  The real part of the lift is the level; its imaginary part is
the winding.  Euler and Weierstrass are two presentations of this one
element.

\subsection{Horizontal: unique orbit--stabilizer factorization}

Every projective arrow has the unique factorization
\[
  f
  =
  \operatorname{orbitRep}(Y)\,
  \operatorname{stab}(f)\,
  \operatorname{orbitRep}(X)^{-1}.
\]
At every instant of every GPV tape it supplies a two-legged commuting
square:
\[
\begin{tikzcd}[column sep=huge,row sep=large]
  \text{input action at }X
    \arrow[r,"D_X(t)"]
    \arrow[d,"{\operatorname{right}(f,t)}"']
  & \text{A-output at }X
    \arrow[d,"{\operatorname{left}(f,t)}"] \\
  \text{input action at }Y
    \arrow[r,"D_Y(t)"']
  & \text{A-output at }Y .
\end{tikzcd}
\]
In multiplicative form,
\[
  \operatorname{left}(f,t)\,D_X(t)
  =
  D_Y(t)\,\operatorname{right}(f,t).
\]
The winding lies in the stabilizer factor.  The orbit factor performs the
horizontal sweep.  Identity and composition of these squares are inherited
from uniqueness of orbit--stabilizer factorization.

\section{The genuine total is a total of the round trip}

The underlying Grothendieck carrier is
\[
  \TA
  =
  \int_{\mathrm{GreatCircle.Base}}\cA.
\]
An object is a projective base point together with the complete normalized
physical/presentation state.  A morphism consists of a base arrow and a
fibre arrow after the corresponding projective reindexing.

\medskip
\noindent
This carrier becomes the genuine total of the A-section only when the input
and output are functors on \emph{total arrows}:
\[
  \In_{\TA},\Out_{\TA}\colon\TA\longrightarrow H_1,
\]
and the octonionic round trip commutes on the total:
\[
\begin{tikzcd}[column sep=huge,row sep=large]
  \TA
    \arrow[r,"\In_{\TA}"]
    \arrow[dr,"\Out_{\TA}"']
  & H_1 \arrow[d,"A_{\OO}"] \\
  & H_1 .
\end{tikzcd}
\qquad
\Out_{\TA}=A_{\OO}\circ\In_{\TA}.
\]
For a total morphism $\langle f,\varphi\rangle$, the fibre morphism
$\varphi$ occurs after $\cA(f)$.  The projective/GPV action square is
precisely the compatibility that makes $\In_{\TA}$ and $\Out_{\TA}$
functorial on this composite arrow.

\paragraph{Why this is the missing gate.}
An object function $\operatorname{Ob}(\TA)\to\operatorname{Ob}(H_1)$ can
store output values, but it cannot express naturality.  A cocone can be
natural only on a functor.  Therefore the output cocone is native to the
construction only after $\Out_{\TA}$ acts on arrows and the triangle above
is proved.

\section{The natural $9+3$ partition}

The twelve readings are not installed as twelve fields and are not
reintroduced after totalization.

\begin{center}
\begin{tabular}{p{0.42\textwidth}p{0.49\textwidth}}
\toprule
\textbf{Nine action readings} & \textbf{Three codomain/output readings} \\
\midrule
Euler level, the complete GPV lift, endpoint real-level conservation,
winding, complex exponential fibre level and height, the octonionic chart,
and orbit--stabilizer transport.
&
The residue-$\CC$ sphere population, its infinitude, and the intrinsic real
coordinate of each realized sphere. \\
\bottomrule
\end{tabular}
\end{center}

The nine readings spread through $\cA$ by the certified action squares.  The
three output readings occur in the functorial output $\Out_{\TA}$.  Their
partition is therefore a consequence of one equivariant total, not a later
conjunction of hypotheses.

\section{The output population and the boundary before 8.3.4}

For every index $n$ and sphere world $I$, the A-action produces a total
object carrying together:
\begin{itemize}[leftmargin=2em]
  \item the normalized octonionic point on the residue-$\CC$ sphere;
  \item its value under $A_{\OO}$;
  \item its intrinsic real register;
  \item the complete Euler--Weierstrass--GPV presentation;
  \item its native orbit--stabilizer branch to the $N$-positioned state.
\end{itemize}

No common value at $N$ is included in the carrier.  Before the next theorem
is applied, the proof has established only the following:
\[
\begin{tikzcd}[column sep=huge]
  \{\text{A-generated residue spheres}\}
    \arrow[r,hook]
  & \TA
    \arrow[r,shift left=0.8ex,"\In_{\TA}"]
    \arrow[r,shift right=0.8ex,"\Out_{\TA}"']
  & H_1 .
\end{tikzcd}
\]
The family is now arrow-natural and output-bearing.  Its common component
and common real centre remain conclusions.

\section{Only now: the categorical readout}

After the total-level triangle is kernel-green, the existing machinery is
applied in this order:
\[
\begin{tikzcd}[column sep=large]
  \TA
    \arrow[r,"\pi_0"]
  & \pi_0(\TA)
    \arrow[r,"\sim" description,
      "\text{Riehl 8.3.4}"]
  & \displaystyle
    \operatorname*{colim}_{X}
      \pi_0\bigl(\cA(X)\bigr)
    \arrow[r,mapsto]
  & \kappa
    \arrow[r,"\operatorname{val}"]
  & c\in\RR .
\end{tikzcd}
\]
Here the component diagram, its cocone, quotient representatives, and
descent are the internal categorical machinery specialized to the exact
total.  They are not predeclared replacements.  The atomic conclusion is
\[
  \exists c\in\RR,\quad
  \forall n\in\mathbb N,\quad
  \operatorname{Re}(A.\operatorname{sphereZero}(n))=c.
\]

\section{Lean obligations and naming key}

\begin{description}[leftmargin=3.3cm,style=nextline]
  \item[\texttt{AsectionSlice}]
    The slice/Moebius projection.  This name is clean and final.
  \item[\texttt{AsectionEquivariant}]
    The actual octonionic round-trip functor $H_1\to H_1$.
  \item[\texttt{AsectionFunctor}]
    The current name of the projective groupoid-valued indexing diagram,
    not the octonionic round trip.  A name such as
    \texttt{AsectionActionDiagram} would expose its role.
  \item[\texttt{TotalA}]
    The current Grothendieck carrier.  It remains provisional until
    total-level naturality is green.
  \item[\texttt{totalInput}/\texttt{totalOutput}]
    Current object functions.  They must be promoted to total input/output
    functors.
\end{description}

\paragraph{Exact Gate 4 certificate.}
\begin{verbatim}
AsectionTotalInput A  : A.TotalA --> H1
AsectionTotalOutput A : A.TotalA --> H1

AsectionTotalInput A ; AsectionEquivariant A =
  AsectionTotalOutput A
\end{verbatim}
The theorem machinery remains gated until these declarations and their
projective/GPV arrow naturality compile.

\paragraph{Arrow-level diagnostic.}
\begin{quote}
Anything that must act on arrows must be checked to actually act on arrows.
\end{quote}
A single endpoint does not retain a continuous lift, and an object function
does not retain a functor.  Both errors erase the arrow dimension while
leaving an object-level term that can still type-check.

\end{document}
